\documentclass[11pt]{article}

\usepackage[letterpaper,margin=0.82in]{geometry}
\usepackage{amsmath}
\usepackage{amssymb}
\usepackage{array}
\usepackage{booktabs}
\usepackage{longtable}
\usepackage{xcolor}
\usepackage{xurl}
\usepackage{hyperref}
\usepackage{enumitem}
\usepackage{titlesec}

\definecolor{DocBlue}{HTML}{1F4E79}
\definecolor{DocRule}{HTML}{D0D7DE}

\hypersetup{
  colorlinks=true,
  linkcolor=DocBlue,
  urlcolor=DocBlue,
  pdfauthor={Crazyflie Simulator},
  pdftitle={Closed-Loop Poles of the Cascaded PID Hover Model}
}

\setlist[itemize]{topsep=0.25em,itemsep=0.12em,parsep=0pt,leftmargin=1.35em}
\renewcommand{\arraystretch}{1.16}
\setlength{\parindent}{0pt}
\setlength{\parskip}{0.42em}
\setlength{\emergencystretch}{3em}
\allowdisplaybreaks

\pagestyle{empty}

\titleformat{\section}
  {\Large\bfseries\color{DocBlue}}
  {\thesection.}{0.55em}{}
\titleformat{\subsection}
  {\large\bfseries\color{DocBlue}}
  {\thesubsection}{0.45em}{}

\DeclareRobustCommand{\code}[1]{\nolinkurl{#1}}
\newcommand{\degc}{d_{\mathrm{eg}}}
\newcommand{\gammar}{\gamma_{\mathrm{rp}}}
\newcommand{\gammay}{\gamma_{\mathrm{yaw}}}

\begin{document}

\begin{center}
{\Huge\bfseries Closed-Loop Poles}\\[0.25em]
{\large Cascaded PID controller + hover-linearized Crazyflie dynamics}
\end{center}


\section{Main Modeling Assumptions}

The closed-loop pole calculation uses the local continuous-time hover model:

\begin{itemize}
  \item level hover, zero yaw, zero angular rates;
  \item small perturbations, no controller saturation or clipping;
  \item quadratic translational drag has zero analytic Jacobian at zero speed;
  \item static wrench-to-PWM-to-thrust conversion is treated as locally identity;
  \item outer-loop integrators use \(K_i^\star=K_i/5\), because the code updates
        the outer loop every five 500 Hz controller ticks.
\end{itemize}

Important numerical constants are

\[
\begin{aligned}
m &= 0.0400,\qquad
g = 9.81,\\
I_x &= I_y = 2.39\times 10^{-5},\qquad
I_z = 3.23\times 10^{-5},\\
c_x &= c_y = 2.0\times 10^{-6},\qquad
c_z = 3.0\times 10^{-6},\\
\gammar &= 4.315969000436706\times 10^{-7},\\
\gammay &= 6.94941672257353\times 10^{-8},\\
\degc &= 180/\pi = 57.29577951308232.
\end{aligned}
\]

Here \(\gammar\) and \(\gammay\) are the code's scale factors from the
firmware-like rate-controller outputs to physical roll/pitch and yaw moments.

\begin{longtable}{>{\raggedright\arraybackslash}p{0.26\linewidth}
                  >{\raggedright\arraybackslash}p{0.28\linewidth}
                  >{\raggedright\arraybackslash}p{0.34\linewidth}}
\toprule
Loop & Code gains & Continuous gains used here \\
\midrule
XY position & \(K_p=1.6,\ K_i=8.0\) & \(K_p=1.6,\ K_{pi}^{\star}=1.6\) \\
XY velocity & \(K_p=22.0,\ K_i=0.5\) & \(K_v=22.0,\ K_{vi}^{\star}=0.1\) \\
Z position & \(K_p=2.0,\ K_i=0.5\) & \(K_{pz}=2.0,\ K_{piz}^{\star}=0.1\) \\
Z velocity & \(K_p=25.0,\ K_i=15.0\) & \(K_{vz}=25.0,\ K_{viz}^{\star}=3.0\) \\
Roll/pitch attitude & \(K_a=6.0,\ K_{ai}=3.0\) & same \\
Roll/pitch rate & \(K_r=200.0,\ K_{ri}=400.0,\ K_{rd}=2.5\) & same \\
Yaw attitude & \(K_a=6.0,\ K_{ai}=1.0,\ K_{ad}=0.35\) & same \\
Yaw rate & \(K_r=120.0,\ K_{ri}=16.7\) & same \\
\bottomrule
\end{longtable}

\section{Hover Plant Used For The Pole Equations}

At level hover, the small-signal plant is

\[
\begin{aligned}
\delta \ddot{x} &= g\,\delta\theta, &
\delta \ddot{y} &= -g\,\delta\phi, &
\delta \ddot{z} &= \frac{1}{m}\delta F,\\
\delta\dot{\phi} &= \delta p, &
\delta\dot{\theta} &= \delta q, &
\delta\dot{\psi} &= \delta r,\\
I_x\delta\dot{p} &= \delta M_x-c_x\delta p, &
I_y\delta\dot{q} &= \delta M_y-c_y\delta q, &
I_z\delta\dot{r} &= \delta M_z-c_z\delta r.
\end{aligned}
\]

At zero yaw, the acceleration-to-command map in the cascaded PID is

\[
\theta_{\mathrm{cmd}}=\frac{a_x}{g},
\qquad
\phi_{\mathrm{cmd}}=-\frac{a_y}{g},
\qquad
\delta F=m\,a_z.
\]

Thus \(x\) and \(y\) have the same characteristic equation: the minus sign in
\(\phi_{\mathrm{cmd}}\) cancels the plant sign \(\ddot{y}=-g\phi\).

\section{Closed-Loop Pole Equations}

\subsection{Roll/Pitch Inner Loop}

For roll or pitch, define the inner attitude transfer denominator
\(D_{\mathrm{att}}(s)\). Let \(I=I_x=I_y\) and \(c=c_x=c_y\).

\[
N_{\mathrm{att}}(s)
=
\gammar\degc\,(K_rs+K_{ri})(K_as+K_{ai})
\]

\[
\boxed{
D_{\mathrm{att}}(s)
=
(Is^2+cs)s^2
+\gammar\degc
\left[
(K_rs+K_{ri})(s^2+K_as+K_{ai})
+K_{rd}s^4
\right].
}
\]

The inner-loop poles are the roots of \(D_{\mathrm{att}}(s)=0\).

\subsection{Horizontal Closed-Loop Poles}

The horizontal position and velocity controllers are

\[
C_p(s)=\frac{K_ps+K_{pi}^{\star}}{s},
\qquad
C_v(s)=\frac{K_vs+K_{vi}^{\star}}{s}.
\]

For either \(x\) or \(y\), the characteristic equation is

\[
\boxed{
D_{xy}(s)
=
D_{\mathrm{att}}(s)s^4
+N_{\mathrm{att}}(s)(K_vs+K_{vi}^{\star})
(s^2+K_ps+K_{pi}^{\star})=0.
}
\]

\subsection{Vertical Closed-Loop Poles}

For the vertical channel,

\[
C_{pz}(s)=\frac{K_{pz}s+K_{piz}^{\star}}{s},
\qquad
C_{vz}(s)=\frac{K_{vz}s+K_{viz}^{\star}}{s}.
\]

The characteristic equation is

\[
\boxed{
D_z(s)
=
s^4
+(K_{vz}s+K_{viz}^{\star})
(s^2+K_{pz}s+K_{piz}^{\star})=0.
}
\]


\subsection{Yaw Closed-Loop Poles}

The yaw attitude loop includes derivative-on-measurement. The characteristic
equation is

\[
\boxed{
\begin{aligned}
D_{\psi}(s)
&=
(I_zs^2+c_zs)s^2\\
&\quad
+\gammay\degc\,(K_rs+K_{ri})
\left[(1+K_{ad})s^2+K_as+K_{ai}\right]=0.
\end{aligned}
}
\]

\subsection{Full Decoupled Pole Set}

Under the hover decoupling assumptions,

\[
\boxed{
\mathcal{P}_{\mathrm{cl}}
=
\mathcal{P}_{xy}
\cup
\mathcal{P}_{xy}
\cup
\mathcal{P}_{z}
\cup
\mathcal{P}_{\psi}.
}
\]

The horizontal pole set appears twice, once for \(x\) and once for \(y\).

\section{Numerical Characteristic Equations And Poles}

\subsection{Horizontal X/Y}

\[
\begin{aligned}
D_{xy}(s)
&=
8.57217021{\times}10^{-5}s^8
+4.94773616{\times}10^{-3}s^7\\
&\quad
+3.95658893{\times}10^{-2}s^6
+7.27023216{\times}10^{-1}s^5\\
&\quad
+2.70927427s^4
+4.32089186s^3\\
&\quad
+3.67547329s^2
+1.06115715s
+0.00474791.
\end{aligned}
\]

\[
\boxed{
\begin{aligned}
\mathcal{P}_{xy}=\{&
-51.744986,\ -2.005099,\ -0.499771,\ -0.00454545,\\
&-0.908355\pm 11.714563j,\\
&-0.823741\pm 1.011742j\}.
\end{aligned}
}
\]

\subsection{Vertical Z}

\[
D_z(s)=s^4+25s^3+53s^2+8.5s+0.3.
\]

\[
\boxed{
\mathcal{P}_{z}
=
\{-22.679597,\ -2.149150,\ -0.119934,\ -0.0513188\}.
}
\]

\subsection{Yaw Psi}

\[
\begin{aligned}
D_{\psi}(s)
&=
0.0000323s^4
+0.000648039042s^3\\
&\quad
+0.00295660812s^2
+0.000876775291s
+0.0000664947655.
\end{aligned}
\]

\[
\boxed{
\mathcal{P}_{\psi}
=
\{-13.365607,\ -6.385026,\ -0.173285,\ -0.139211\}.
}
\]

\subsection{Main Conclusion}

For the 12-state hover-linearized plant without motor lag, all closed-loop
poles are in the open left-half plane. The slowest pole is

\[
\boxed{s=-0.00454545\ {\rm s}^{-1}},
\]

from the horizontal cascaded-integral dynamics. The full decoupled pole set
contains \(8+8+4+4=24\) poles because it combines \(x\), \(y\), \(z\), and yaw
channels.

\section{Complete Rotor Failure Case}
\label{sec:rotor-failure}

\subsection{Rotor-Failure Model}

Let the commanded rotor force vector be

\[
f_{\mathrm{cmd}}=
\begin{bmatrix}f_1&f_2&f_3&f_4\end{bmatrix}^{T}.
\]

The healthy mixer maps rotor forces to body wrench by

\[
w =
\begin{bmatrix}F&M_x&M_y&M_z\end{bmatrix}^{T}
=G f,
\]

\[
G =
\begin{bmatrix}
1&1&1&1\\
-a&-a&a&a\\
-a&a&a&-a\\
-k&k&-k&k
\end{bmatrix}.
\]

If rotor \(i\) completely fails, define

\[
E_i=\operatorname{diag}(1,\ldots,0,\ldots,1),
\]

where the zero is in the failed rotor's diagonal entry. The realized wrench for
a commanded wrench is then

\[
\boxed{
\delta w_{\mathrm{real}}
=
H_i\,\delta w_{\mathrm{cmd}},
\qquad
H_i=G E_i G^{-1}.
}
\]

The closed-loop state used for the pole calculation is

\[
\xi =
\begin{bmatrix}
\delta x_{\mathrm{plant}}^{T} &
\delta x_{\mathrm{PID\ integrators}}^{T}
\end{bmatrix}^{T}
\in \mathbb{R}^{24}.
\]

Because the roll/pitch rate loop uses derivative-on-measurement
\(-K_{rd}\dot{p}\) and \(-K_{rd}\dot{q}\), the failed-actuator closed-loop model
is most cleanly written as a descriptor system:

\[
\boxed{
\mathcal{E}_i\dot{\xi}=\mathcal{A}_i\xi.
}
\]

The post-failure pole set is therefore

\[
\boxed{
\mathcal{P}_{\mathrm{fail},i}
=
\left\{s\in\mathbb{C}:
\det(s\mathcal{E}_i-\mathcal{A}_i)=0
\right\}.
}
\]

For the default repo fault, rotor 1 fails. The corresponding commanded-wrench
to realized-wrench map is

\[
H_1 =
\begin{bmatrix}
0.750000 & 7.071068 & 7.071068 & 43.915210\\
0.008839 & 0.750000 & -0.250000 & -1.552637\\
0.008839 & -0.250000 & 0.750000 & -1.552637\\
0.001423 & -0.040254 & -0.040254 & 0.750000
\end{bmatrix}.
\]

\subsection{General Case}

The exact \(H_1\) matrix above has off-diagonal terms, so the failed-rotor
system is not truly separable into independent \(x\), \(y\), \(z\), and yaw
channels. However, if one wants the same channel-wise format as the healthy
case, a useful approximation is to keep only the same-channel effectiveness
terms:

\[
H_{1,\mathrm{diag}}
=
\operatorname{diag}(0.75,\ 0.75,\ 0.75,\ 0.75).
\]

This means each nominal channel keeps the same structure as before, but with
75 percent control effectiveness. The resulting approximate channel poles are:

\[
\boxed{
\begin{aligned}
\mathcal{P}_{x,\mathrm{fail}} \approx
\mathcal{P}_{y,\mathrm{fail}} \approx \{&
-47.080680,\ -2.005563,\ -0.499771,\ -0.00454545,\\
&-0.789809\pm11.754564j,\\
&-0.823755\pm1.011763j\}.
\end{aligned}
}
\]

\[
\boxed{
\mathcal{P}_{z,\mathrm{fail}} \approx
\{-16.341340,\ -2.237428,\ -0.119912,\ -0.0513195\}.
}
\]

\[
\boxed{
\mathcal{P}_{\psi,\mathrm{fail}} \approx
\{-7.379053\pm3.092900j,\ -0.173235,\ -0.139226\}.
}
\]

This channel-wise view is useful for intuition: losing one rotor reduces
effective authority and moves the fast poles. It does not replace the exact
coupled failed-rotor pole set below.

\subsection{Numerical Poles For Rotor 1 Complete Failure}

For baseline cascaded PID with rotor 1 completely failed, no motor lag, and no
saturation, the generalized eigenvalues of
\((\mathcal{A}_1,\mathcal{E}_1)\) are

\[
\boxed{
\begin{aligned}
\mathcal{P}_{\mathrm{fail},1}=\{&
-51.744986,\ -44.759105,\ -18.588507,\ -3.596782,\\
&-2.006860,\ -2.005099,\\
&-0.908355\pm 11.714563j,\\
&-0.824472\pm 1.009904j,\\
&-0.823741\pm 1.011742j,\\
&-0.616459\pm 11.154612j,\\
&-0.499771,\ -0.499752,\\
&-0.142809\pm0.014255j,\\
&-0.010123,\ -0.00454545,\ -0.00454545,\\
&-6.72622\times10^{-5},\\
&3.36299\times10^{-5}\pm5.80496\times10^{-5}j
\}.
\end{aligned}
}
\]

The last complex pair:

\[
\boxed{
s=3.36299\times10^{-5}\pm5.80496\times10^{-5}j.
}
\]

lies in the right-half plane.

\begin{longtable}{>{\raggedright\arraybackslash}p{0.18\linewidth}
                  >{\raggedright\arraybackslash}p{0.24\linewidth}
                  >{\raggedright\arraybackslash}p{0.42\linewidth}}
\toprule
Failed rotor & Largest real part & Low-frequency unstable mode \\
\midrule
\endfirsthead
\toprule
Failed rotor & Largest real part & Low-frequency unstable mode \\
\midrule
\endhead
1 & \(3.36299\times10^{-5}\) &
\(3.36299\times10^{-5}\pm5.80496\times10^{-5}j\) \\
2 & \(3.36299\times10^{-5}\) &
\(3.36299\times10^{-5}\pm5.80496\times10^{-5}j\) \\
3 & \(9.73686\times10^{-5}\) &
\(9.73686\times10^{-5}\) \\
4 & \(3.36299\times10^{-5}\) &
\(3.36299\times10^{-5}\pm5.80496\times10^{-5}j\) \\
\bottomrule
\end{longtable}

Thus any single complete rotor failure makes the baseline full-yaw cascaded PID
lose local asymptotic stability in this small-signal model.

\subsection{If Yaw Is Removed From The Model}

If yaw is not considered, remove the yaw state, yaw controller, and commanded
yaw moment, and keep the exact failed-rotor coupling from commanded
\([F,M_x,M_y]^T\) to realized \([F,M_x,M_y]^T\):

\[
H_{1,xyz} =
\begin{bmatrix}
0.750000 & 7.071068 & 7.071068\\
0.008839 & 0.750000 & -0.250000\\
0.008839 & -0.250000 & 0.750000
\end{bmatrix}.
\]

The resulting 20-state \(x/y/z\) closed-loop model has poles

\[
\boxed{
\begin{aligned}
\mathcal{P}_{xyz,\mathrm{fail},1}=\{&
-51.744986,\ -43.636141,\ -9.289775,\ -2.480127,\\
&-2.006220,\ -2.005099,\\
&-0.908355\pm11.714563j,\\
&-0.823827\pm1.011870j,\\
&-0.823741\pm1.011742j,\\
&-0.321372\pm11.630657j,\\
&-0.499771,\ -0.499771,\\
&-0.119868,\ -0.0513210,\\
&-0.00454545,\ -0.00454545
\}.
\end{aligned}
}
\]

% All of these poles are in the left-half plane. Thus, in this local model, the
% instability caused by complete rotor failure is tied to trying to keep the full
% yaw channel in the loop. Once yaw is removed, the \(x/y/z\) subsystem can be
% locally stable, although the real vehicle can still spin in yaw and the
% allocator can still saturate.

% \section{What Motor Lag Changes}
% \label{sec:motor-lag}

% \subsection{Does The Current Simulation Include Motor Lag?}

% Yes. In \code{model/quadcopter.py}, the runtime path is

% \[
% \text{PWM command}
% \rightarrow
% \text{desired rotor thrust}
% \rightarrow
% \text{first-order thrust lag}
% \rightarrow
% \text{realized wrench}
% \rightarrow
% \text{state dynamics}.
% \]

% The first-order lag is asymmetric:

% \[
% \tau =
% \begin{cases}
% 0.025\ {\rm s}, & \text{if desired rotor thrust is increasing},\\
% 0.040\ {\rm s}, & \text{if desired rotor thrust is decreasing}.
% \end{cases}
% \]

% This means the runtime simulator is piecewise/nonlinear even near hover. There
% is no unique exact LTI motor-lag pole set unless we freeze \(\tau\).

% \subsection{Frozen-Tau Pole Equations}

% If one freezes the actuator as

% \[
% P_m(s)=\frac{1}{\tau s+1},
% \]

% then motor lag modifies the characteristic denominators as follows:

% \[
% \boxed{
% D_{\mathrm{att},\tau}(s)
% =
% (\tau s+1)(Is^2+cs)s^2
% +\gammar\degc
% \left[
% (K_rs+K_{ri})(s^2+K_as+K_{ai})
% +K_{rd}s^4
% \right].
% }
% \]

% \[
% \boxed{
% D_{xy,\tau}(s)
% =
% D_{\mathrm{att},\tau}(s)s^4
% +N_{\mathrm{att}}(s)(K_vs+K_{vi}^{\star})
% (s^2+K_ps+K_{pi}^{\star}).
% }
% \]

% \[
% \boxed{
% D_{z,\tau}(s)
% =
% (\tau s+1)s^4
% +(K_{vz}s+K_{viz}^{\star})
% (s^2+K_{pz}s+K_{piz}^{\star}).
% }
% \]

% \[
% \boxed{
% \begin{aligned}
% D_{\psi,\tau}(s)
% &=
% (\tau s+1)(I_zs^2+c_zs)s^2\\
% &\quad
% +\gammay\degc\,(K_rs+K_{ri})
% \left[(1+K_{ad})s^2+K_as+K_{ai}\right].
% \end{aligned}
% }
% \]

% \subsection{Frozen-Tau Numerical Poles}

% \begin{longtable}{>{\raggedright\arraybackslash}p{0.18\linewidth}
%                   >{\raggedright\arraybackslash}p{0.74\linewidth}}
% \toprule
% Case & Poles \\
% \midrule
% \(\tau=0.025,\ x/y\) &
% \(-68.740586\pm50.651764j,\ -2.005029,\ -0.956482\pm11.817861j,\ -0.823742\pm1.011739j,\ -0.499771,\ -0.00454545\) \\
% \(\tau=0.025,\ z\) &
% \(-18.846238\pm23.610824j,\ -2.136270,\ -0.119934,\ -0.0513188\) \\
% \(\tau=0.025,\ \psi\) &
% \(-17.020042\pm17.463897j,\ -5.740299,\ -0.173285,\ -0.139211\) \\
% \midrule
% \(\tau=0.040,\ x/y\) &
% \(-41.809302\pm52.511293j,\ -2.004987,\ -0.987670\pm11.883022j,\ -0.823743\pm1.011738j,\ -0.499771,\ -0.00454545\) \\
% \(\tau=0.040,\ z\) &
% \(-11.349993\pm21.061735j,\ -2.128761,\ -0.119934,\ -0.0513188\) \\
% \(\tau=0.040,\ \psi\) &
% \(-9.638974\pm17.170464j,\ -5.502435,\ -0.173286,\ -0.139211\) \\
% \bottomrule
% \end{longtable}

% The practical interpretation is: motor lag mainly adds or moves fast actuator
% poles. It barely changes the slowest closed-loop poles in this tuning. For a
% simulation-consistent small-signal discussion, report the frozen-\(\tau\) poles
% as a piecewise approximation; for the repo's nominal 12-state hover model,
% report the healthy-case poles in Section 4.

\end{document}
